% 分野別類題: 行列の対角化・固有値・固有ベクトル
% 2x2・3x3 対称行列の直交対角化と、対角化を利用した A^n の計算 3 問。
% コンパイル: TEXINPUTS="../../../00_common:$TEXINPUTS" latexmk -xelatex problems.tex
\documentclass[a4paper,11pt]{article}
\usepackage{preamble}

\begin{document}

\kamokumei{類題演習 --- 行列の対角化・固有値・固有ベクトル}

\shijibun{次の各問について、行列 $A$ の固有値と固有ベクトルを求め、$A$ を対角化せよ（対称行列については直交行列 $P$ による直交対角化 $P^{-1}AP = P^{\mathrm{T}}AP = D$ まで示すこと）。導出過程を明示せよ。}

\shomon{問1.}{次の対称行列 $A$ の固有値・固有ベクトルを求め、直交行列 $P$ を用いて直交対角化 $P^{\mathrm{T}}AP = D$ を行え。}
\[
A = \begin{pmatrix} 3 & 1 \\ 1 & 3 \end{pmatrix}
\]

\shomon{問2.}{次の対称行列 $A$ の固有値・固有ベクトルを求め、直交行列 $P$ を用いて直交対角化 $P^{\mathrm{T}}AP = D$ を行え。}
\[
A = \begin{pmatrix} 2 & 1 & 1 \\ 1 & 2 & 1 \\ 1 & 1 & 2 \end{pmatrix}
\]

\shomon{問3.}{次の行列 $A$ を対角化し、それを用いて $A^{n}$（$n$ は自然数）を求めよ。}
\[
A = \begin{pmatrix} 1 & 2 \\ 2 & 1 \end{pmatrix}
\]

\end{document}
